\section{The Geometry of Roots}

So far we have treated roots as elements of the dual vector space \(\mathfrak{h}^{\vee}\) of the Cartan subalgebra \(\mathfrak{h}\), which is a complex vector space. In this section we will see how they have a geometric meaning in real Euclidean space, which makes them (almost) visualizable.

\begin{theorem}
    This theorem gives two statements which let us embed the roots in a real Euclidean space.
    \begin{enumerate}
        \item Let \(\mathfrak{h}_\mathbb{R} \subset \mathfrak{h}\) be the real vector space spanned by \(H_\alpha \in \mathcal{R}\) for some \(\alpha \in \mathcal{R}\). Then \(\mathfrak{h} = \mathfrak{h}_\mathbb{R} \oplus i\mathfrak{h}_\mathbb{R} \cong (\mathfrak{h}_\mathbb{R})_\mathbb{C}\). The restriction of the Killing form to \(\mathfrak{h}_\mathbb{R}\) is \(\mathbb{R}\)-valued and positive definite.
        \item Let \(\mathfrak{h}^{\vee}_\mathbb{R} \subset \mathfrak{h}^{\vee}\) be the real vector space spanned by \(\alpha \in \mathcal{R} \subset \mathfrak{h}^{\vee}\). Then \(\mathfrak{h}^{\vee} = (\mathfrak{h}^{\vee}_\mathbb{R})_\mathbb{C}\). Also, \(\mathfrak{h}^{\vee}_\mathbb{R} = \{\lambda \in \mathfrak{h}^{\vee} | \lambda(H) \in \mathbb{R} \forall H \in \mathfrak{h}_\mathbb{R}\} = (\mathfrak{h}_\mathbb{R})^{\vee}\).
    \end{enumerate}
\end{theorem}

[...]

[...]

\begin{example}[\(\mathfrak{g} = \mathfrak{sl}(2, \mathbb{C})\)]
    [...]

    [...]

\end{example}

\begin{example}[\(\mathfrak{g} = \mathfrak{sl}(3, \mathbb{C})= \mathfrak{su}(3)_\mathbb{C}\)]

    [...]

    For an alternative approach, we quote a fact without proof, which is that for a simple

        [...]

    Thus we have the following Euclidean representation of the roots (up to an overall rescaling of their lengths):

    [\textit{drawing}]

\end{example}

The \(\mathfrak{sl}(3, \mathbb{C})\) example shows that roots can be written as integer linear combinations of a subset of them. To make this precise, we have the following definitions.